\section{Related work}

The target of this paper is the homogeneous conditional log-odds coefficient \leanref{sym:\theta}{\ensuremath{\theta(P)}} in \cref{obj:def:logistic-class}. The statistical experiment combines this pointwise homogeneity restriction with a uniform covariate design and prescribed bounds on the nuisance logits and their Hölder seminorms, as specified in \cref{obj:def:model}. Precision is measured by expected interval length on effect-radius subclasses, while coverage is required over the ambient model; \cref{obj:def:radius-model,obj:def:honest-procedures,obj:def:length-objective} formalize these distinct roles. These choices connect the analysis to logistic coefficient estimation, estimation of nonlinear functionals under low smoothness, and honest confidence intervals.

The closest logistic coefficient results develop robustness and efficiency within partially linear logistic models. \citet[Propositions 1--3]{Tan2019LogisticDR} characterizes nuisance-orthogonal estimating functions and obtains unbiasedness when either the baseline-logit model or the conditional exposure-mean model given covariates and a fixed outcome value is correct. The latter nuisance is an outcome-conditioned exposure mean. When both nuisance models are correct, the efficiency comparison identifies an optimal weight within the proposed estimating-function class. The accompanying near-null comparison explains why a simpler weight approaches that optimal choice as the coefficient approaches zero. This supplies a useful precedent for effect-dependent precision: proximity to the null changes the efficiency comparison within a regular coefficient-estimation problem. Here, proximity to the null indexes subclasses on which the length of an ambiently honest interval is evaluated.
% [Tan (2019), Propositions 1--3 and efficiency discussion](https://arxiv.org/pdf/1901.09138)

Direct odds-ratio antecedents make the nuisance comparison especially clear. \citet{TchetgenTchetgen2010} constructs doubly robust estimators in a semiparametric odds-ratio model by pairing alternative baseline conditional-density specifications. For binary exposure, \citet{TchetgenTchetgen2013} gives a closed-form adjusted odds-ratio estimator with the corresponding double-robustness interpretation. Those papers provide estimation and large-sample inference for an odds-ratio parameter under their nuisance specifications. The present result instead fixes a homogeneous scalar conditional log odds, requires coverage over the full prescribed Hölder model, evaluates length over every effect-radius subclass, and proves same-model converses as well as attainment.

Regular asymptotic inference for the logistic coefficient is developed by \citet[arXiv version 1, Assumptions REG and ML1, Theorem 2]{LiuZhangZhou2020LogisticDML}. Their regularity conditions impose compact parameter and exposure domains together with nondegenerate score sensitivity and variance. For sample size \leanref{sym:n}{\ensuremath{n}}, Assumption ML1 requires uniform consistency of the baseline-logit and outcome-conditioned exposure-mean estimators, with both \(L^2\) errors of order \(o_P(n^{-1/4})\). Under these conditions, the coefficient estimator admits an asymptotically linear representation and a normal limit at the parametric rate. The present expected-length criterion evaluates precision over the prescribed Hölder classes, including their low-smoothness regimes, subject to the coverage requirement in \cref{obj:ass:global-coverage}.
% [Liu, Zhang, and Zhou (2020), arXiv version 1](https://arxiv.org/pdf/2009.14461v1)

The causal interpretation of the estimand also aligns with DINA, the difference in natural parameters. For binary responses, \citet[Section 2.1.1, equation (5)]{GaoHastie2025DINA} defines DINA as the conditional log-odds contrast between potential outcomes. A constant treatment-effect basis specializes this formulation to a homogeneous effect. Their finite-dimensional coefficient analysis assumes a correctly represented treatment-effect function, bounded covariates and nuisance functions, and a score derivative bounded away from singularity. Under those conditions, nuisance estimation errors enter the coefficient bound quadratically, alongside a parametric sampling term \citep[Section 3.3, Proposition 1]{GaoHastie2025DINA}. This comparison separates the shared conditional log-odds target from the precision criterion: their coefficient bound controls estimation error, while the present objective optimizes expected interval length under ambient coverage.
% [Gao and Hastie (2025), Sections 2.1.1 and 3.3](https://academic.oup.com/biometrics/article/81/4/ujaf162/8403946)

The observable numerator in \cref{obj:lem:observable-odds} connects the construction to expected conditional covariance estimation. In the known-design setting, \citet[version 5, Theorem 1 and Corollary 1]{McGrath2026} gives projection bias of order \(k^{-s/d}\) and variance bounded by a constant times \(n^{-1}+k/n^2\), with \(s\) the sum of nuisance smoothness exponents and \(d\) the covariate dimension; the resulting resolution choices attain the covariance minimax rate. Related results extend to local averaging methods: \citet[arXiv version 3, Assumption 4 and Theorem 2]{McCleanBalakrishnanKennedyWasserman2024DCDR} obtains expected absolute error \(n^{-2s/(2s+d)}\) when \(s\le d/2\) and root-\(n\) inference when \(s>d/2\), using independent nuisance-training samples, a known smooth covariate density bounded away from zero, and suitable undersmoothing. Below the regular threshold, additional undersmoothing yields a slower normal limit under boundedness and continuity of conditional variances \citep[arXiv version 3, Theorem 3]{McCleanBalakrishnanKennedyWasserman2024DCDR}. These results distinguish rate-optimal estimation from normal inference with negligible bias. The present construction combines those projection ingredients for two observable coordinates with ambient honesty, all-radius length evaluation, and a homogeneous-logistic converse, as specified in \cref{obj:def:coordinate-estimates,obj:def:upper-handle}.
% [McGrath, version 5, Theorem 1 and Corollary 1](https://arxiv.org/pdf/2212.14857v5)
% [McClean et al., arXiv version 3, Assumption 4 and Theorems 2--3](https://arxiv.org/html/2403.15175v3)

Higher-order inference for implicitly defined parameters provides a further comparison. \citet[Section 3.3, Assumptions 1 and 3--9, Theorems 1--4]{ZhangLiuZhang2026GeneralTreatmentHOIF} analyzes solutions to corrected moment equations under identification, a nonsingular moment derivative, boundedness and overlap, complexity and continuity conditions on the moment class, and localized, nondegenerate sieve bases. Their oracle error bound combines a stochastic term with the product of two residual projection errors. Feasible estimation adds a product involving nuisance errors, the initial parameter estimate, and the error in estimating the projection Gram matrix. Normal inference requires the locally uniform projection bias to vanish after normalization, an additional continuity-rate condition, and, for the feasible estimator, the stated Gram-matrix error condition. The normalization accommodates sieve dimensions both below and above sample size, with dimension growing more slowly than \(n^2\). This decomposition clarifies the shared role of projection in reducing nuisance bias. In the prescribed logistic experiment, the coordinate identity permits an explicit interval construction whose bias and variance bounds enter the finite-sample coverage argument in \cref{obj:thm:finite-honest-upper}.
% [Zhang, Liu, and Zhang (2026), Section 3.3](https://arxiv.org/html/2606.01706v1)

Converse results must be compared at the level of both estimand and statistical class. \citet[Appendix B]{JinMackeySyrgkanis2025HardNormal} studies structure-agnostic bounds for binary treatment in a partially linear mean model, making explicit the roles of overlap and feasible nuisance perturbations. For average conditional log-odds differences, \citet[version 2, Theorem 7.7]{JinSyrgkanis2025GeneralLinearLower} obtains an absolute-error lower bound containing a squared regression-error term, a mixed nuisance-error term, and the parametric sampling term. That result uses unrestricted log-odds contrasts in nuisance-error neighborhoods, under its density, overlap, nondegeneracy, and nonzero-curvature conditions. Applying a converse to the present coefficient requires alternatives that satisfy the homogeneous logistic relation throughout the covariate domain, together with the prescribed envelopes and smoothness restrictions. The calibrated families and original-record mixture comparisons in \cref{obj:def:continuous-calibrated-priors,obj:lem:calibrated-support,obj:lem:original-record-mixtures} supply that comparison within the stated model.
% [Jin, Mackey, and Syrgkanis (2025), Appendix B](https://arxiv.org/html/2507.02275v1)
% [Jin and Syrgkanis, version 2, Theorem 7.7](https://arxiv.org/html/2512.17341v2)

Finally, the distinction between the coverage class and the length-evaluation class belongs to the honest-interval literature. \citet{Low1997} relates expected length at a distribution to coverage over a family of distributions. In Gaussian experiments for linear functionals, \citet[Section 2, Theorem 1]{CaiLow2004ConfidenceAdaptation} expresses a corresponding subclass expected-length lower bound through a between-class modulus of continuity. The present criterion adopts this separation for a homogeneous logistic experiment: the public smoothness exponents specify the ambient class, and effect radii specify nested subclasses for evaluating length. The contribution of \cref{obj:thm:full-honest-frontier-answer} is the joint all-radius characterization under its stated conditions, combining an attaining ambiently honest procedure with matching expected-length lower comparisons, including evaluation at the null.
% [Cai and Low (2004), Section 2, including the discussion of Low (1997)](https://arxiv.org/pdf/math/0503662)

The research provenance also includes two internal supporting problems: the known-fair-radius development \texttt{stat\_dina\_logodds\_honest\_radius} and the fixed-exponent \((0.3,0.1)\) development \texttt{stat\_logodds\_anisotropic\_radius\_frontier}. They supplied antecedent calibration and lower-bound architectures; in particular, the parametric null-floor reduction is credited to the known-fair development. The present theorem establishes the complete public-exponent domain and all-radius profile. The neighboring internal problem \texttt{stat\_logodds\_roughdesign\_radius\_frontier} studies a different experiment with fixed exponents and an unknown bounded design density, whereas the current design is known and uniform.
